\begin{lemma}
\label{editorial-source-lemma-lift-idempotents}
Let $R$ be a ring. Let $I \subset R$ be a locally nilpotent ideal.
Then $R \to R/I$ induces a bijection on idempotents.
\end{lemma}

\begin{proof}[First proof of Lemma \ref{editorial-source-lemma-lift-idempotents}]
As $I$ is locally nilpotent it is contained in every prime ideal.
Hence $\Spec(R/I) = V(I) = \Spec(R)$. Hence the
lemma follows from Lemma \ref{algebra-lemma-disjoint-decomposition}.
\end{proof}

\begin{proof}[Second proof of Lemma \ref{editorial-source-lemma-lift-idempotents}]
Suppose $\overline{e} \in R/I$ is an idempotent.
We have to lift $\overline{e}$ to an idempotent of $R$.

\medskip\noindent
First, choose any lift $f \in R$ of $\overline{e}$, and set
$x = f^2 - f$. Then, $x \in I$, so $x$ is nilpotent (since $I$
is locally nilpotent). Let now $J$ be the ideal of $R$ generated
by $x$. Then, $J$ is nilpotent (not just locally nilpotent),
since it is generated by the nilpotent $x$.

\medskip\noindent
Now, assume that we have found a lift $e \in R$ of $\overline{e}$
such that $e^2 - e \in J^k$ for some $k \geq 1$.
Let $e' = e - (2e - 1)(e^2 - e) = 3e^2 - 2e^3$, which is another
lift of $\overline{e}$ (since the idempotency of $\overline{e}$
yields $e^2 - e \in I$). Then
$$
(e')^2 - e' = (4e^2 - 4e - 3)(e^2 - e)^2 \in J^{2k}
$$
because expansion of either side gives
$4e^6-12e^5+9e^4+2e^3-3e^2$.
In particular every coefficient and sign in the original update is retained.

\medskip\noindent
Choose $N\geq1$ with $x^N=0$, so $J^N=0$.
Starting with $e_0=f$, define for every $j\geq0$
$$
e_{j+1}=e_j-(2e_j-1)(e_j^2-e_j)=3e_j^2-2e_j^3.
$$
The identity just proved gives inductively
$e_j^2-e_j\in J^{2^j}$, and every $e_j$ still maps to
$\overline e$ since each update differs by an element of $J$.
Any integer $j$ with $2^j\geq N$ therefore gives
$e_j^2-e_j=0$. This constructs the required lift after finitely
many explicitly specified updates.

\medskip\noindent
We thus have seen that if $\overline{e} \in R/I$ is any
idempotent, then there exists a lift of $\overline{e}$
which is an idempotent of $R$.
It remains to prove that this lift is unique. Indeed, let
$e_1$ and $e_2$ be two such lifts. We
need to show that $e_1 = e_2$.

\medskip\noindent
By definition of $e_1$ and $e_2$, we have $e_1 \equiv e_2
\mod I$, and both $e_1$ and $e_2$ are idempotent. From
$e_1 \equiv e_2 \mod I$, we see that $e_1 - e_2 \in I$,
so that $e_1 - e_2$ is nilpotent (since $I$ is locally nilpotent).
A straightforward
computation (using the idempotency of $e_1$ and $e_2$)
reveals that $(e_1 - e_2)^3 = e_1 - e_2$. Using this and
induction, we obtain $(e_1 - e_2)^k = e_1 - e_2$ for any
positive odd integer $k$. Since all high enough $k$ satisfy
$(e_1 - e_2)^k = 0$ (since $e_1 - e_2$ is nilpotent),
this shows $e_1 - e_2 = 0$, so that $e_1 = e_2$, which
completes our proof.
\end{proof}